\documentclass[11pt]{article}

\usepackage[utf8]{inputenc}
\usepackage[T1]{fontenc}
\usepackage{geometry}
\usepackage{amsmath}
\usepackage{amssymb}
\usepackage{booktabs}
\usepackage{hyperref}
\usepackage{xurl}
\usepackage{xcolor}
\usepackage{listings}
\usepackage{enumitem}
\usepackage{parskip}
\usepackage{graphicx}
\graphicspath{{../images/}}

\geometry{margin=1in}

\lstset{
  language=Java,
  basicstyle=\ttfamily\small,
  keywordstyle=\color{blue},
  commentstyle=\color{gray},
  stringstyle=\color{teal},
  showstringspaces=false,
  columns=fullflexible,
  frame=single,
  framerule=0.2pt,
  breaklines=true
}

\setlist[itemize]{topsep=0.3em,itemsep=0.2em}
\setlist[enumerate]{topsep=0.3em,itemsep=0.2em}

% Simple per-section source line.
\newcommand{\notesources}[1]{\par\noindent{\small\color{gray}\textit{Sources:} \sloppy #1\par}}

% Unit-wise source strings for this subject (used by slides/notes).
% Path is relative to the lecture `latex/` folder (the compilation CWD).
% OOPS with Java (CSEG1044) — unit-wise sources
%
% These are short, student-facing reference strings intended to be used with:
%   - \slidesources{...}  (in beamer slides)
%   - \notesources{...}   (in lecture notes)
%
% Style inspiration: other subjects in My-Subjects/ use semicolon-separated sources
% and include an "accessed" date for web documentation.

\newcommand{\OOPSUnitISources}{Schildt, \textit{Java: A Beginner's Guide} (10e), McGraw-Hill (Java fundamentals + OOP basics).; Oracle Java Tutorials: Getting Started + Language Basics + Classes and Objects (accessed \today).; Java SE API docs (e.g., \texttt{String}, \texttt{Scanner}) (accessed \today).}

\newcommand{\OOPSUnitIISources}{Schildt, \textit{Java: The Complete Reference} (12e), McGraw-Hill (classes, inheritance, packages, interfaces).; Bloch, \textit{Effective Java} (3e), Addison-Wesley, 2018 (design choices; interfaces vs abstract classes).; Oracle Java Tutorials: Classes and Objects + Interfaces + Packages (accessed \today).; Java SE API docs (e.g., \texttt{Comparable}, \texttt{Runnable}) (accessed \today).}

\newcommand{\OOPSUnitIIISources}{Schildt, \textit{Java: The Complete Reference} (12e) (nested/inner classes, exceptions, threads, I/O).; Oracle Java Tutorials: Nested Classes + Exceptions + Concurrency + Basic I/O (accessed \today).; Java SE API docs (e.g., \texttt{Thread}, \texttt{AutoCloseable}) (accessed \today).}

\newcommand{\OOPSUnitIVSources}{Schildt, \textit{Java: The Complete Reference} (12e) (generics, lambdas, Swing, JDBC).; Bloch, \textit{Effective Java} (3e) (generics + lambdas).; Oracle Java Tutorials: Generics + Lambda Expressions + Swing + JDBC Basics (accessed \today).; Java SE API docs (e.g., \texttt{java.util.function}, \texttt{javax.swing}, \texttt{java.sql}) (accessed \today).}

\newcommand{\OOPSUnitVSources}{Schildt, \textit{Java: The Complete Reference} (12e) (Collections Framework + wrapper classes).; Oracle Java docs: Collections Framework overview (accessed \today).; Java SE API docs (\texttt{java.util} collections; wrapper classes in \texttt{java.lang}) (accessed \today).}

\newcommand{\OOPSUnitVISources}{Git SCM: \textit{Pro Git} (2e) (version control workflow), accessed \today.; GitHub Docs (repositories, issues, pull requests), accessed \today.; Oracle Java Tutorials: Swing + JDBC (accessed \today).; JUnit 5 User Guide (accessed \today).; Apache Maven Project: Getting Started (accessed \today).; Gradle User Manual: Getting Started (accessed \today).; UML class diagrams: UML 2.x overview/spec (accessed \today).}

\newcommand{\OOPSMiscWhyInterfacesSources}{Schildt, \textit{Java: The Complete Reference} (12e) (interfaces).; Bloch, \textit{Effective Java} (3e) (prefer interfaces where appropriate).; Oracle Java Tutorials: Interfaces (accessed \today).; Java SE API docs (e.g., \texttt{Comparable}, \texttt{Runnable}) (accessed \today).}



\title{Object Oriented Programming (CSEG1044)\\Miscellaneous Notes\\Why Do We Need Interfaces?}
\author{Tofik Ali}
\date{\today}

\begin{document}
\maketitle

\section*{The core idea (1 minute)}
An \textbf{interface} lets you describe \textit{what an object can do} (its behavior) without committing to \textit{how it does it} (its implementation).

\paragraph{One-line answer.} We need interfaces so our code can depend on \textbf{contracts}, not concrete classes --- which makes systems easier to \textbf{change}, \textbf{extend}, and \textbf{test}.

\section*{What goes wrong without interfaces (tight coupling)}
Consider a checkout flow.

\paragraph{Naive approach.} The service creates a specific payment class internally:
\begin{lstlisting}
class CardPayment {
    void pay(int amount) { /* call bank */ }
}

class CheckoutService {
    void checkout(int amount) {
        CardPayment payment = new CardPayment(); // fixed choice
        payment.pay(amount);
    }
}
\end{lstlisting}

\paragraph{Problems.}
\begin{itemize}
  \item If tomorrow you add UPI / Wallet payments, you must edit \texttt{CheckoutService}.
  \item You cannot easily test \texttt{CheckoutService} without calling the real bank API.
  \item Your high-level business logic is tightly tied to low-level implementation details.
\end{itemize}

\section*{Interfaces solve this by introducing a contract}
Define a contract for payment:
\begin{lstlisting}
interface PaymentMethod {
    void pay(int amount);
}
\end{lstlisting}

Provide multiple implementations:
\begin{lstlisting}
class CardPayment implements PaymentMethod {
    public void pay(int amount) { /* call bank */ }
}

class UpiPayment implements PaymentMethod {
    public void pay(int amount) { /* call UPI */ }
}
\end{lstlisting}

Now the checkout logic depends only on the contract:
\begin{lstlisting}
class CheckoutService {
    private final PaymentMethod payment;

    CheckoutService(PaymentMethod payment) {
        this.payment = payment;
    }

    void checkout(int amount) {
        payment.pay(amount);
    }
}
\end{lstlisting}

\paragraph{What changed?}
\begin{itemize}
  \item \texttt{CheckoutService} no longer cares \textit{which} payment method is used.
  \item You can add \texttt{WalletPayment} later without changing \texttt{CheckoutService}.
  \item This is the idea behind ``program to an interface, not an implementation''.
\end{itemize}

\section*{Use case 1: Plug-in behavior (Strategy pattern)}
Interfaces are a clean way to plug in behavior.

\paragraph{Example.} Same checkout flow; different pricing rules:
\begin{lstlisting}
interface DiscountPolicy {
    int discount(int total);
}

class NoDiscount implements DiscountPolicy {
    public int discount(int total) { return 0; }
}

class StudentDiscount implements DiscountPolicy {
    public int discount(int total) { return (int)(0.10 * total); }
}

class Billing {
    private final DiscountPolicy policy;
    Billing(DiscountPolicy policy) { this.policy = policy; }
    int finalAmount(int total) { return total - policy.discount(total); }
}
\end{lstlisting}

\paragraph{Takeaway.} You can swap policies at runtime without rewriting the billing code.

\section*{Use case 2: Testability (fake/mock implementations)}
Interfaces make unit testing practical because you can replace a real dependency with a fake.

\paragraph{Example.} A notification service:
\begin{lstlisting}
interface Notifier {
    void send(String to, String message);
}

class EmailNotifier implements Notifier {
    public void send(String to, String message) { /* send email */ }
}

class OrderService {
    private final Notifier notifier;
    OrderService(Notifier notifier) { this.notifier = notifier; }

    void placeOrder(String userEmail) {
        // business logic...
        notifier.send(userEmail, "Order placed");
    }
}
\end{lstlisting}

\paragraph{Testing without sending real emails.}
\begin{lstlisting}
class FakeNotifier implements Notifier {
    String lastTo;
    String lastMessage;
    public void send(String to, String message) {
        lastTo = to;
        lastMessage = message;
    }
}
\end{lstlisting}

\paragraph{Takeaway.} The interface lets you test \texttt{OrderService} fast, safely, and deterministically.

\section*{Use case 3: Decouple storage (repositories/data access)}
Many applications use an interface to separate business logic from data storage.

\begin{lstlisting}
class User { String id; String name; }

interface UserRepository {
    User findById(String id);
    void save(User u);
}

class InMemoryUserRepository implements UserRepository {
    // can be used for demos/tests
    public User findById(String id) { return null; }
    public void save(User u) { }
}

class SqlUserRepository implements UserRepository {
    // can be used in production
    public User findById(String id) { return null; }
    public void save(User u) { }
}
\end{lstlisting}

\paragraph{Takeaway.} Business code can use \texttt{UserRepository} without caring about SQL, files, or APIs.

\section*{Why interfaces are especially useful in Java}
\begin{itemize}
  \item \textbf{Multiple roles:} a class can implement multiple interfaces (Java does not allow multiple inheritance of classes).
  \item \textbf{Framework callbacks:} Java uses interfaces heavily for ``I call you'' designs.
\end{itemize}

\paragraph{Examples from Java.}
\begin{itemize}
  \item Collections: \texttt{List}, \texttt{Set}, \texttt{Map}, \texttt{Iterable}, \texttt{Iterator}
  \item Ordering: \texttt{Comparable}, \texttt{Comparator}
  \item Concurrency: \texttt{Runnable}, \texttt{Callable}
  \item I/O and lifecycle: \texttt{Closeable}, \texttt{AutoCloseable}
\end{itemize}

\section*{Interface vs abstract class (simple rule)}
\begin{itemize}
  \item Prefer an \textbf{interface} when you want to define a capability/role and expect multiple unrelated implementations.
  \item Prefer an \textbf{abstract class} when you need shared state and shared code across closely related classes.
  \item In modern Java, interfaces can also contain \texttt{default} methods --- use them carefully (they are still a contract-first tool).
\end{itemize}

\section*{Mini exercises (in class or homework)}
\begin{enumerate}
  \item Design a small ``export'' feature using an \texttt{Exporter} interface.\\
        Provide two implementations: \texttt{CsvExporter} and \texttt{JsonExporter}.
  \item Write a tiny unit test using a \texttt{FakeNotifier} to verify that an order confirmation message is sent.
  \item Refactor any one program you wrote earlier: replace direct use of a concrete class with an interface + two implementations.
\end{enumerate}
\end{document}

\notesources{\OOPSMiscWhyInterfacesSources}
